\scanpage{52}
\mspage{8}

On a ainsi la notion de \uline{graphe pondéré} \deleted{\unsure{enrichi}} \uline{à repères} $\unsure{\Gamma|r}$

\margin{Groupoïde $(\uline{\mathrm{Gr.pondrep}})$\\ foncteur\\ $S_{*}\unsure{\mathscr T|\mathbb T}{\unsure{\mathbb D|D}}_{\mathrm{top}}^{\pm}\longrightarrow(\mathrm{Gr\,pond\,rep})$}

\diagram{\begin{tikzcd}[ampersand replacement=\&,column sep=small]
{\vec R_l^{+}\ \deleted{\mathbb D_\infty}\subset\underset{\sigma_0,\sigma_1}{\vec R}}\arrow[r] \&
\vec A\arrow[loop right,"\sigma"]\arrow[d,"o"']\&\\
\&S\arrow[r,"\gamma"]\&\mathbb N
\end{tikzcd}}

qui sera importante sans doute dans l'explicitation de l'opération de ``vrillement''. Il est clair en tous cas que toute description-traduction algébrico-combinatoire de $\unsure{S\mathscr T|\mathscr T}\mathbb D_{\mathrm{top}}^{\pm}$ devra permettre, pour chaque objet du \deleted{la catégorie} \inserted{groupoïde-}paradigme, de réconstituer (à isom. unique près) le \uline{graphe pondéré à repères} correspondant.

\inserted{--- plus précisément, il faut réconstituer le foncteur canonique $S_{*}\mathscr T\mathbb D_{\mathrm{top}}^{\pm}\longrightarrow(\mathrm{Grpondrep.})$ (et non seul\textsuperscript{mt} le foncteur à valeurs dans $(\mathrm{Grbipond})$)}

\uline{Opérations essentielles}

A) Opérations ``d'oubli'' ou de ``\uline{gommage}''.

\textcircled{1} \uline{Gommage de repères} \quad \unsure{A'|À} un objet $(X,\Sigma,R)$, muni en plus d'une partie $R'\subset R$, il associe $(X,\Sigma,R\setminus R')$ (gommage de $R'$).
% archive sheet number: 51

